% Pillar-7: E1 open-backward-pass gradient validation (toy scale, directional)
\section{E1: An Open-Backward-Pass Check of T3}
\label{sec:p7-e1}
Theorems T1 and T2 are checkable from reward telemetry alone, but T3---that
gradient magnitude tracks $p(1-p)$---requires seeing the gradients, which
the closed stack used for most of the benchmark does not expose. Experiment
E1 (W\&B project \texttt{zvf-colab-experiments}, run family
\texttt{E1\_grad\_signal}) therefore reruns the loop in an open trainer with
a fully instrumented backward pass: a 0.5B-parameter model on a synthetic
arithmetic task with binary verifiable rewards, chosen so that per-step
success probability sweeps a wide range of $p$ within a single run.

\paragraph{Status: anecdote, not evidence.} An earlier revision reported
$\mathrm{corr}(\|\nabla\|,\hat p(1-\hat p))=+0.71$ from this run. We demote it
to an anecdote and do not use it to support T3, for three reasons. (i) It
cannot be recomputed: only the scalar survives in the run summary; the
per-step gradient norms and group success rates were not retained. (ii) It was
computed over roughly 30 per-step rows from a single run, which are serially
dependent, so no meaningful interval or test is available. (iii) It is largely
mechanical: with group-centered advantages, a step whose groups are all
degenerate has zero advantage and hence a near-zero policy-gradient term, so a
positive association between $\|\nabla\|$ and within-group contrast is expected
from the loss definition alone (Appendix~\ref{sec:p7-appendix}) and is not an
independent test of the binomial model. T3 therefore stands as a modelling
assumption; the test that would support it is D3 in
Section~\ref{sec:p7-limitations}.
